\begin{example}
    Let's solve the system
    \[\begin{cases}
        2x+4y =3 \\
        -x+y =2
    \end{cases}\]
    using one of the two methods we have, substitution.

    We first solve by substitution -- the best approach when dealing with only two equations and two unknowns.

    To solve a system by substitution, you will solve one equation for one of the variables and plug this expression into the other equation.

    Let's solve the bottom equation for \(y\) and plug this representation of \(y\) into the top equation.

    \[y = x+2 \rightarrow 2x+4\left(x+2\right) = 3.\]

    From this new equation in the single variable \(x,\) we can determine the \(x\) coordinate of the solution to the system of equations.

    \[2x+4\left(x+2\right) = 3 \rightarrow x = -\frac{5}{6}.\]

    We can plug this \(x\) into either of the other two equations to get the \(y\) coordinate of the solution.

    \[-\frac{5}{6} + 2 = \frac{7}{6}.\]

    The only solution to this system of equations is \(\left(-\frac{5}{6},\frac{7}{6}\right).\)
\end{example}

\begin{example}
    Let's solve the system
    \[\begin{cases}
        x+3y=9 \\
        \frac{1}{3}x+2y =2
    \end{cases}\]
    using the other of the two methods we have, elimination. The goal of this method is to ``eliminate'' variables from one of the equations. There are two things that we can do to the system that will not affect the system of equations. We can multiply an equation by any real number. This does not affect the set of solutions for the single equation and hence does not affect the set of solutions for the entire solutions. We can also add these scaled equations together. While this creates a new equations which has a different set of solutions, the set of solutions to both equations is preserved since these are the points at which we have equality.

    We will refer to \(x+3y=6\) as ``top'' and \(\frac{1}{3}x+y =2\) as ``bottom.'' We may eliminate the \(x\) from bottom by replacing bottom with \(-\frac{1}{3}top + bottom.\)

    \[\begin{matrix}
        &-x/3-y & = -3 \\
        + & x/3+2y & = 2 \\
        \hline & 0+y & =-1
    \end{matrix}\]

    Our new system.

    \[\begin{cases}
        x+3y=9 \\
        y =-1
    \end{cases}\]

    Now we can complete this problem by substituting \(y=-1\) into top; however, we continue with the elimination process.

    We now wish to eliminate \(y\) from the top equation. We accomplished this by replacing top with \(top-3bottom.\)

    \[\begin{matrix}
        &x+3y & =9 \\
        - &3y & =-3 \\
        \hline & x & = 12
    \end{matrix}\]

    Thus our solution is \(\left(12,-1\right).\)
\end{example}

We can solve a system with three variables and three unknowns with elimination, as well.